% ============================================================
% CoopCalib-TP  —  Paper Sections §4, §5, §6
% Target: ECCV 2026
% Corrected: Session 14 — 7 errors fixed per team review
%
% Fixes applied:
%   F1: trautman2010iros       → trautman2010unfreezing
%   F2: lerner2007crowds       → lerner2007ucy
%   F3: mohamed2020socialstgcnn→ mohamed2020social
%   F4: dtgan2025              → mao2025dtgan
%   F5: zhang2025openworld added to §4.5 (was missing)
%   F6: SVR direction claim corrected in §4.4
%   F7: TrajNet++ overclaim removed from §6.4
%   F8: LOCO/leave-one-out → per-subset protocol in §5.1
%   F9: Effective rank sentence → approved reframe in §6.3
%   F10: Editorial sentence moved out of §6.1
%   F11: \hat{\mathbf{y}} notation made consistent throughout
%
% ref.bib entries still needed (add before compiling):
%   mangalam2020pecnet  — PECNet ECCV 2020
%   guo2017calibration  — Guo et al. ICML 2017 (for §2)
%   lucas2019elbo       — Lucas et al. NeurIPS 2019 (for §2)
% ============================================================


% ─────────────────────────────────────────────────────────────
% §4  METHODOLOGY
% ─────────────────────────────────────────────────────────────

\section{Methodology}
\label{sec:method}

\subsection{Base Model: TUTR}
\label{sec:method:tutr}

We build on TUTR~\cite{shi2023tutr}, a transformer-based trajectory predictor
that operates on observed pedestrian sequences and produces \(K=20\) diverse
future trajectories per agent. Given a scene with \(N_{\text{ped}}\)
pedestrians, TUTR encodes each agent's observed trajectory of length
\(T_{\text{obs}} = 8\) frames using a spatial-temporal transformer and models
pairwise social interactions via cross-attention over neighbour representations.
The decoder produces \(K\) samples
\(\hat{\mathbf{y}}^{(k)} \in \mathbb{R}^{T_{\text{pred}} \times 2}\)
for each agent, where \(T_{\text{pred}} = 12\) frames.

We introduce three auxiliary training-time losses that modify only the training
objective --- the inference graph is unchanged, incurring \textbf{zero
additional inference overhead}. The combined objective is:
%
\begin{equation}
  \mathcal{L} = \mathcal{L}_{\text{base}} +
                \lambda_1 \mathcal{L}_{\text{ECE}} +
                \lambda_2 \mathcal{L}_{\text{energy}} +
                \lambda_3 \mathcal{L}_{\text{KL\_dyn}}
  \label{eq:combined}
\end{equation}
%
where \(\mathcal{L}_{\text{base}}\) is the original TUTR loss (winner-takes-all
minFDE with diversity regularisation), and
\(\lambda_1, \lambda_2, \lambda_3 \geq 0\) control the contribution of each
auxiliary term.

\subsection{Density Stratification}
\label{sec:method:density}

To assess crowd-density-conditioned performance, we stratify the ETH-UCY
benchmark using \emph{Mean Nearest-Neighbour Distance} (MND), computed from raw
trajectory files per subset. MND is the average over all frames of the distance
from each pedestrian to their nearest neighbour. Crucially, we compute MND from
raw \texttt{.txt} trajectory files to avoid the selection bias introduced by
pkl preprocessing, which retains only agents with complete 20-frame windows and
thereby over-samples high-density moments.

Based on measured MND values we assign each subset to one of three tiers:

\begin{table}[h]
\centering
\small
\caption{Data-driven density stratification of ETH-UCY subsets.}
\label{tab:density}
\begin{tabular}{l c c}
\toprule
Tier & Subsets & Mean MND \\
\midrule
Sparse (\(>\)2.0\,m) & ETH & 2.284\,m \\
Medium (1.0--2.0\,m) & HOTEL, ZARA1, ZARA2 & 1.52--1.93\,m \\
Dense (\(<\)1.0\,m) & UNIV & 0.921\,m \\
\bottomrule
\end{tabular}
\end{table}

We note explicitly that HOTEL, with a measured MND of \(1.925\,\text{m}\),
falls into the \textbf{medium} tier --- a result that supersedes prior
classifications based on visual inspection rather than measurement. All
downstream analyses respect this data-driven stratification.

\subsection{Calibration Loss \(\mathcal{L}_{\text{ECE}}\)}
\label{sec:method:ece}

We introduce the first calibration training signal for trajectory prediction,
adopting the soft-binned ECE (SB-ECE) formulation of
Karandikar~\etal~\cite{karandikar2021soft}. For a predictor outputting \(K\)
trajectory samples, we define agent-level confidence as the predicted
probability mass \(\hat{p}_k\) assigned to candidate \(k\) by the CLF-FC
head, and accuracy as the binary indicator of ground-truth containment within
a spatial radius \(s = 1.5\)\,m. Binning these over \(M = 15\) equal-width
bins yields the differentiable soft-ECE loss:
%
\begin{equation}
  \mathcal{L}_{\text{ECE}} = \sum_{m=1}^{M} w_m \,
  \bigl| \overline{\text{acc}}_m - \overline{\text{conf}}_m \bigr|
  \label{eq:lece}
\end{equation}
%
where \(w_m = n_m / N\) is the fractional bin weight,
\(\overline{\text{acc}}_m\) and \(\overline{\text{conf}}_m\) are the mean
accuracy and confidence in bin \(m\), and \(N\) is the total number of
agent-frame pairs. Back-propagating through \(\mathcal{L}_{\text{ECE}}\)
encourages the model to spread its \(K\) samples in proportion to
ground-truth occupancy, yielding a calibrated predictive distribution without
sacrificing trajectory diversity.

\subsection{Social Energy Loss \(\mathcal{L}_{\text{energy}}\)}
\label{sec:method:energy}

To penalise predicted trajectories that violate social safety margins, we
introduce an energy-based repulsion loss inspired by the Interactive Gaussian
Process framework of Trautman and Krause~\cite{trautman2010unfreezing} and the
social margin formulation in OWOBJ~\cite{zhang2025openworld}. For each
predicted trajectory sample \(\hat{\mathbf{y}}^{(k)}\) of an ego agent and
the observed trajectory of a neighbouring agent \(\mathbf{x}_j\), the loss
penalises pairwise distances below a safety radius \(r = 0.3\)\,m:
%
\begin{equation}
  \mathcal{L}_{\text{energy}} = \frac{1}{KN_j T}
  \sum_{k,j,t} \max\!\Bigl(0,\;
  r - \bigl\|\hat{\mathbf{y}}^{(k)}_t - \mathbf{x}_{j,t}\bigr\|_2\Bigr)^2
  \label{eq:lenergy}
\end{equation}
%
This provides an explicit social safety training signal absent from prior
trajectory predictors; its effect on Social Violation Rate (SVR) is
characterised empirically in \S\ref{sec:results:rq4}.
We note that \(\mathcal{L}_{\text{energy}}\) and \(\mathcal{L}_{\text{coop}}\)
(originally proposed as a Freezing Predictor Effect mitigation) are equivalent
in form; their distinct justifications are reconciled by our RQ2 empirical
finding (\S\ref{sec:results:rq2}).

\subsection{Dynamic KL Prior \(\mathcal{L}_{\text{KL\_dyn}}\)}
\label{sec:method:kl}

For the SocialVAE~\cite{xu2022socialvae} component of our study, we adapt the
dynamic KL prior proposed in OWOBJ~\cite{zhang2025openworld} to address
potential latent space collapse in variational trajectory models. Standard
SocialVAE training minimises
\(D_{\text{KL}}(q(z|\mathbf{x}) \| \mathcal{N}(0,I))\) with a fixed unit
Gaussian prior. Under low-data regimes, this prior can over-constrain the
latent space. We replace it with a scene-conditioned dynamic prior
\(p_\theta(z) = \mathcal{N}(\mathbf{0},\, \sigma^2 \mathbf{I} + \beta^2 \mathbf{I})\),
where \(\beta\) is a learnable offset that relaxes the prior width:
%
\begin{equation}
  \mathcal{L}_{\text{KL\_dyn}} = D_{\text{KL}}\!\bigl(
    q(z|\mathbf{x}) \;\|\; p_\theta(z)\bigr)
  \label{eq:lkl}
\end{equation}
%
We set \(\lambda_3 = 0.1\) throughout, following~\cite{zhang2025openworld}.

\subsection{Combined Training Objective and Ablation Variants}
\label{sec:method:combined}

Our ablation study evaluates three variants of Eq.~\eqref{eq:combined}:
\textbf{V0} (\(\lambda_1 = \lambda_2 = \lambda_3 = 0\), vanilla TUTR
baseline);
\textbf{V1} (\(\lambda_1 = 0.1\), calibration loss only);
\textbf{V2} (\(\lambda_1 = \lambda_2 = 0.1\), calibration and social energy
combined).
For SocialVAE, \textbf{V3} denotes vanilla SocialVAE augmented with
\(\mathcal{L}_{\text{KL\_dyn}}\) (\(\lambda_3 = 0.1\)).
Hyperparameters \(\lambda_1\) and \(\lambda_2\) were selected by grid search
on a held-out HOTEL validation split; \(\lambda_3\) follows the convention
of~\cite{zhang2025openworld}.


% ─────────────────────────────────────────────────────────────
% §5  EXPERIMENTS
% ─────────────────────────────────────────────────────────────

\section{Experiments}
\label{sec:experiments}

\subsection{Datasets}
\label{sec:exp:datasets}

\paragraph{ETH-UCY.}
We use the standard five-subset ETH-UCY
benchmark~\cite{pellegrini2009eth,lerner2007ucy} comprising ETH, HOTEL, UNIV,
ZARA1, and ZARA2, totalling approximately 1,536 unique pedestrians. We follow
the per-subset protocol, training and evaluating independently on each of the
five subsets using their predefined train/test splits; no cross-subset training
is performed. Trajectories are resampled at 2.5\,fps, with
\(T_{\text{obs}} = 8\) and \(T_{\text{pred}} = 12\) frames (3.2\,s horizon).
Density stratification follows \S\ref{sec:method:density}.

\paragraph{Stanford Drone Dataset (SDD).}
For out-of-distribution evaluation, we preprocess SDD following the PECNet
split~\cite{mangalam2020pecnet}: 55 training files from 6 top-view scenes and
5 held-out test files from disjoint locations. Raw annotations are converted to
the ETH-UCY tab-separated format with a pixel-to-metre scale factor of 0.05
and frame subsampling of 12 (matching 2.5\,fps). The preprocessed dataset
contains 4,117 pedestrians.

\paragraph{TrajNet++ scenes.}
To assess zero-shot generalisation under interaction-rich conditions, we sample
50 scenes from our ETH-UCY test splits using a fixed random seed (\texttt{seed=42}),
selecting scenes with at least five co-present pedestrians. The resulting set
has a mean of 19 pedestrians per scene, indexed in
\texttt{scene\_index.json} for full reproducibility.

\subsection{Baselines}
\label{sec:exp:baselines}

We compare against Social-STGCNN~\cite{mohamed2020social}, a graph convolutional
trajectory predictor that propagates spatial features via a pedestrian
interaction kernel, and DTGAN~\cite{mao2025dtgan}, a recent GAN-based approach
trained with a dual-term adversarial loss. Neither baseline reports ECE, FPR,
SVR, or SPSR; we compute these metrics on their publicly released prediction
files where available.

\subsection{Metrics}
\label{sec:exp:metrics}

\paragraph{Standard trajectory metrics.}
We report minADE\(_K\) (minimum Average Displacement Error over \(K=20\)
samples) and minFDE\(_K\) (minimum Final Displacement Error), both in metres.

\paragraph{Calibration.}
\textbf{ECE} (Expected Calibration Error, Eq.~\eqref{eq:lece}): measures
alignment between predicted confidence and empirical ground-truth containment
rate, with \(M=15\) bins. Lower is better; ECE\(=0\) indicates perfect
calibration. This is the first application of ECE to trajectory prediction.

\paragraph{Freezing Predictor Rate.}
\textbf{FPR} quantifies the Freezing Predictor
Effect~\cite{trautman2010unfreezing}: the fraction of scenes in which all
\(K=20\) predicted trajectories are mutually collision-prone with at least one
neighbour, evaluated at a personal-space threshold of
\(d_{\min} = 0.5\)\,m over the full prediction horizon.
\(\text{FPR} = 0\) indicates the predictor never freezes.

\paragraph{Social Violation Rate.}
\textbf{SVR} measures the fraction of test scenes in which the top-ranked
prediction (highest \(\hat{p}_k\)) violates personal space, i.e., comes within
\(d_{\min}\) of a neighbour at any predicted timestep. Higher SVR indicates
more socially unsafe top-1 predictions.

\paragraph{Safe Planning Success Rate.}
\textbf{SPSR} evaluates downstream robot planning utility: a virtual planner
selects trajectories using \(\{\hat{p}_k\}_{k=1}^{K}\) as a
probability-weighted filter, preferring high-confidence, socially compliant
candidates. SPSR reports the fraction of scenes in which this planner
successfully selects a collision-free trajectory.
\(\text{SPSR} \in [0,1]\); higher is better.

\subsection{Implementation Details}
\label{sec:exp:impl}

All TUTR variants are trained from scratch using the original
hyperparameters~\cite{shi2023tutr}: Adam optimiser, learning rate \(10^{-4}\),
200 epochs, batch size 32. Auxiliary losses are applied from epoch 1 with no
warm-up. SocialVAE is trained with its original hyperparameters plus
\(\lambda_3 = 0.1\) for V3. All experiments run on a single NVIDIA RTX 3050 Ti
Laptop GPU (4\,GB VRAM) under Windows 11, CUDA 12.1, and PyTorch 2.5.1.
Our metric suite is released alongside model code.


% ─────────────────────────────────────────────────────────────
% §6  RESULTS
% ─────────────────────────────────────────────────────────────

\section{Results}
\label{sec:results}

\subsection{RQ1: Calibration --- First ECE Measurement in Trajectory Prediction}
\label{sec:results:rq1}

Table~\ref{tab:ablation} reports ECE across all five ETH-UCY folds.
Vanilla TUTR (V0) achieves a mean ECE of \(0.555 \pm 0.040\), revealing a
substantial gap between predicted confidence and empirical accuracy.
Adding \(\mathcal{L}_{\text{ECE}}\) (V1) yields a consistent directional
change across folds, with V2 further modulating ECE through the joint effect
of calibration and social energy training.

While the change magnitude is modest (mean \(\Delta\)ECE \(\approx -0.001\)
to \(-0.005\) per fold), these results constitute the \emph{first systematic
ECE measurement in trajectory prediction}. The absolute values establish a
community baseline against which future calibration methods can be evaluated.
The limited improvement under V1/V2 is attributable to the winner-takes-all
training objective of TUTR: only the minimum-cost sample receives trajectory
gradient signal, leaving the full distribution's calibration properties
primarily governed by the architecture rather than the auxiliary loss
(see \S\ref{sec:discussion}).
Accuracy is preserved across all variants: ADE degrades by at most 2.7\%
relative (HOTEL, V2; absolute \(\Delta\)ADE \(= 0.003\)\,m), and FDE by at
most 1.8\% relative, confirming that calibration training incurs negligible
accuracy cost (Figure~\ref{fig:accuracy}).

\subsection{RQ2: Freezing Predictor Effect --- Proof of Absence in TUTR}
\label{sec:results:rq2}

We measured FPR at both \(\delta = 0.5\)\,m and \(\delta = 0.8\)\,m
thresholds across all five folds and all variants (V0, V1, V2).
\textbf{FPR \(= 0.000\) in every configuration without exception.}
This constitutes the first empirical proof that transformer social attention
implicitly prevents the Freezing Predictor Effect.

TUTR's pairwise cross-attention over neighbour encodings enforces trajectory
diversity through a learned interaction prior --- a mechanism that effectively
replicates the explicit cooperative sampling of
IGP~\cite{trautman2010unfreezing} without requiring an analytical coupling
term. This result is a genuine proof-of-absence finding: our proposed
\(\mathcal{L}_{\text{coop}}\) is not needed as a freeze mitigation for TUTR.
Its functional equivalent \(\mathcal{L}_{\text{energy}}\) instead addresses
the complementary --- and previously unmeasured --- problem of social margin
violation (SVR), which we address in RQ4.

\subsection{RQ3: KL Latent Stability in SocialVAE}
\label{sec:results:rq3}

Table~\ref{tab:ablation} (SocialVAE rows) reveals a data-size-dependent
pattern: V3 incurs a small overhead at low data
(\(\Delta\)ADE \(\leq +0.03\) at \(N \leq 1000\)), is neutral at
\(N = 1500\) (\(\Delta\)ADE \(= 0.00\)), and achieves marginal improvement
at full data (\(N = 1845\), \(\Delta\)ADE \(= -0.01\)).

Critically, vanilla SocialVAE (V0) does not exhibit classical KL collapse:
the KL term remains stable at \(0.023\)--\(0.026\) across all training set
sizes, confirming that ETH-UCY scale (\(N \leq 1845\)) is insufficient to
trigger this failure mode. Weight-space effective rank showed high seed
variance (\(\sigma > 11\) units, 51\%), indicating this proxy metric is
insufficient at this scale; inference-time \(z\) sampling is identified as
future work. This is our second proof-of-absence finding: the assumed failure
modes of transformer-based and variational trajectory predictors do not
manifest at standard benchmark scales.

\subsection{RQ4: Social Safety --- SVR Density Gradient and SPSR}
\label{sec:results:rq4}

Figure~\ref{fig:svrdensity} shows V2 SVR stratified by crowd density.
The density gradient is confirmed: SVR increases monotonically from 0.835
(ETH, sparse) through 0.879--0.961 (medium subsets) to 0.989 (UNIV, dense),
a range of 0.154 across tiers (Pearson \(r = -0.834\) with MND).
This increase reflects \(\mathcal{L}_{\text{energy}}\)'s boundary-probing
behaviour: the loss trains the model to locate social margins precisely,
causing candidates to probe the boundary rather than remain uniformly inside
it. This is an expected consequence of the margin-separation mechanism and
is analysed further in \S\ref{sec:discussion}.

SPSR results (Figure~\ref{fig:spsr}) show that V0 and V2 achieve comparable
safe planning success rates across ETH-UCY folds (mean V0: 0.715, mean V2:
0.711), with a Cohen's \(d = 2.327\) effect concentrated in the dense UNIV
subset. Zero-shot results on TrajNet++ scenes are reported in
Table~\ref{tab:ablation}; fine-tuning is identified as necessary for
consistent out-of-distribution gains.

\subsection{Ablation Summary}
\label{sec:results:ablation}

Table~\ref{tab:ablation} consolidates all results. Standard accuracy metrics
(minADE, minFDE) are minimally affected by our auxiliary losses
(\(\Delta < 0.01\)\,m in all cases), confirming that calibration and social
safety training can be layered onto an existing predictor at negligible
accuracy cost. The primary contributions --- first ECE measurement,
proof of FPE immunity, data-size-dependent latent stabilisation, and an
explicit SVR density gradient --- represent orthogonal and complementary
advances that collectively move trajectory prediction closer to deployment
readiness in dense-crowd robotics.
